\ticket{27}{Сетевое планирование проекта. Примеры создания сетевых графиков и назначения ресурсов}

\begin{examplan}
    \item Дать определение сетевого планирования.
    \item Описать работы, события, зависимости и сетевые графики.
    \item Раскрыть CPM, PERT, критический путь и расчет сроков.
    \item Объяснить назначение и выравнивание ресурсов.
\end{examplan}

\subsection*{Короткая версия для письменного ответа}

\textbf{Сетевое планирование} --- метод планирования и контроля проекта через граф работ и зависимостей. Он помогает определить длительность проекта, критический путь, резервы времени и влияние ограниченных ресурсов.

\subsection*{Основные понятия}

\begin{description}
    \item[Работа] задача проекта, требующая времени и ресурсов.
    \item[Событие] момент начала или окончания работы, не имеющий длительности.
    \item[Сетевой график] граф работ и логических зависимостей.
    \item[Зависимость] ограничение порядка выполнения работ.
    \item[Путь] последовательность связанных работ от начала к завершению.
    \item[Критический путь] самый длинный путь проекта; его длительность задает минимальный срок проекта.
    \item[Резерв времени] допустимая задержка работы без сдвига срока проекта.
\end{description}

Типы зависимостей: Finish-to-Start, Start-to-Start, Finish-to-Finish, Start-to-Finish. Наиболее распространена Finish-to-Start.

\subsection*{Методы графиков}

\begin{itemize}
    \item \textbf{AOA/ADM}: работы изображаются дугами, события --- узлами. Иногда нужны фиктивные работы нулевой длительности.
    \item \textbf{AON/PDM}: работы изображаются узлами, зависимости --- стрелками. Этот способ чаще используется в современном ПО управления проектами.
\end{itemize}

AOA явно показывает события, но сложнее при изменениях и требует фиктивных работ. AON обычно нагляднее и проще для сложных проектов.

\subsection*{CPM и PERT}

\textbf{CPM} --- метод критического пути. Использует детерминированные длительности работ и вычисляет ранние/поздние сроки и резервы.

\textbf{PERT} применяется при неопределенности длительностей и использует три оценки: оптимистическую, наиболее вероятную и пессимистическую.

\subsection*{Расчет временных параметров}

Для каждой работы:
\begin{itemize}
    \item $ES$ --- раннее начало;
    \item $EF=ES+Dur$ --- раннее окончание;
    \item $LF$ --- позднее окончание;
    \item $LS=LF-Dur$ --- позднее начало;
    \item $TF=LS-ES=LF-EF$ --- полный резерв.
\end{itemize}

Расчет:
\begin{enumerate}
    \item \textbf{Прямой проход}: от начала к концу, считаем $ES$ и $EF$. Для работы $ES$ равен максимуму $EF$ предшественников.
    \item \textbf{Обратный проход}: от конца к началу, считаем $LF$ и $LS$. Для работы $LF$ равен минимуму $LS$ последователей.
    \item Работы с $TF=0$ образуют критический путь.
\end{enumerate}

Любая задержка критической работы задерживает весь проект, если не принять компенсирующие меры.

\subsection*{Диаграмма Ганта}

Диаграмма Ганта показывает работы как полосы на временной шкале. Она удобна для визуализации календарного расписания, параллельных работ, длительностей и фактического прогресса. Сетевой график лучше показывает логику зависимостей, а диаграмма Ганта --- календарное представление.

\subsection*{Назначение ресурсов}

После построения графика назначаются ресурсы: люди, оборудование, материалы, финансы.

Процесс:
\begin{enumerate}
    \item определить ресурсы для каждой работы;
    \item оценить доступность ресурсов;
    \item назначить ресурсы на работы;
    \item построить загрузку ресурсов;
    \item выявить перегрузки и недогрузки;
    \item выполнить выравнивание или сглаживание ресурсов.
\end{enumerate}

\textbf{Resource leveling} может сдвигать работы и даже увеличивать длительность проекта, если ресурсов не хватает. \textbf{Resource smoothing} оптимизирует загрузку в пределах уже заданных сроков, используя резервы некритических работ.

Способы выравнивания: сдвиг некритических работ в пределах резерва, разделение задач, привлечение дополнительных ресурсов, изменение длительности или состава работ.

\begin{important}
    \item Сетевое планирование показывает зависимости между работами.
    \item Критический путь определяет минимальную длительность проекта.
    \item Резерв времени показывает, насколько работу можно задержать без сдвига проекта.
    \item CPM использует прямой и обратный проходы для расчета сроков.
    \item Назначение ресурсов может изменить реальное расписание даже при неизменной логике зависимостей.
\end{important}
